\chapter{Projected Impact}
\label{chap:projected_impact}

\section{Accomplishments and Benefits}
\label{sec:accomplishments}

\begin{enumerate}
    \item \textbf{Consistent Multi-Seed Accuracy \& Low Wrong-Link Rate:} Across 5 random seeds on $12,000$ locked UA-DETRAC test events, GapLink-Lite achieves $\mathbf{93.64\% \pm 0.14\%}$ Top-1 accuracy (vs.\ $80.29\%$ for C-3) while cutting WLR by $6.5\times$ ($19.24\% \to \mathbf{2.98\%}$ 5-seed mean, $3.01\%$ Seed~1).
    \item \textbf{Online Tracking Gains with Standard MOT Metrics:} Across $40$ test sequences ($1,922$ occlusion intervals up to 90 frames, $a_{\max}=95$), GapLink-Lite reaches $\mathbf{96.20}\%$ IRR / $\mathbf{99.11}\%$ IDF1 (full occlusion) and $\mathbf{95.99}\%$ IRR / $\mathbf{99.02}\%$ IDF1 (partial occlusion with zero-leakage visible-crop histograms).
    \item \textbf{Compact Footprint \& Real-Time Live-Detector VMS Deployment:} With $327,979$ parameters ($1.3\text{ MB}$), GapLink-Lite trains in $\sim 100\text{ s}$ per seed and runs end-to-end with live YOLOv8n at $\mathbf{34.5}$--$\mathbf{79.7\text{ FPS}}$ on a laptop RTX 3050 6GB GPU. Across $40$ raw test videos ($3,103$ natural YOLOv8n occlusions), \textbf{GapLink-VMS} reaches $\mathbf{87.30\%}$ natural recovery while cutting forbidden-route false alarms by $3.8\times$ ($49.22\% \to \mathbf{12.88\%}$) and improving occluded anti-loitering dwell F1 by $+17.87\text{ pp}$ over C-3.
\end{enumerate}

\section{Future Improvements}
\label{sec:future_improvements}

\begin{itemize}
    \item \textbf{Detector-Aware Fine-Tuning:} While zero-shot transfer from clean trajectories to live YOLOv8n already cuts natural wrong-link rate by $3.7\times$ ($37.78\% \to 10.20\%$), fine-tuning directly on noisy live-detector box jitter and class-flip artifacts can further suppress residual false associations.
    \item \textbf{Multi-Modal Mixture Predictor:} Extending the single 2D Gaussian predictor to a $K=2$ or $K=3$ Gaussian Mixture Model (GMM) would better model multi-branch intersection turns during long occlusions.
    \item \textbf{Moving-Camera Ego-Motion:} Supporting dashboard cameras will require homography/optical-flow compensation prior to kinematic normalization.
\end{itemize}
